//323111047


\section{Introducción}

Planteamiento del problema. El reto técnico de este proyecto consiste en desarrollar una aplicación de línea de comandos en lenguaje C que simule un sistema de carrito de compras con gestión dinámica de memoria. El problema central radica en solicitar al usuario el tamaño del carrito y permitir la captura de datos mediante estructuras que incluyan el nombre, identificador y precio de cada producto. Para resolver esto, es indispensable administrar correctamente los punteros para el almacenamiento de los elementos y asegurar que el sistema sea capaz de procesar la lista completa al finalizar la ejecución para generar un ticket de compra con el cálculo exacto del precio total de todos los artículos ingresados.

\vspace{0.4cm}

Motivación. La necesidad de dar solución a este problema se basa en la importancia de diseñar software escalable que se adapte a las necesidades del usuario en tiempo real. En el desarrollo de sistemas transaccionales y de comercio electrónico, es imposible predecir con exactitud cuántos artículos seleccionará un cliente, por lo que el uso eficiente de la memoria dinámica es vital para evitar el desperdicio de recursos o la limitación de la capacidad del sistema. Este proyecto permite al estudiante aplicar los conocimientos del primer parcial en un escenario práctico que imita la lógica de funcionamiento de aplicaciones comerciales reales y la optimización de recursos computacionales.

\vspace{0.4cm}

Objetivos. El propósito fundamental de este proyecto es lograr la implementación funcional de un sistema de carrito de compras que utilice memoria dinámica para el almacenamiento de estructuras de productos. Se busca que el alumno demuestre dominio en la creación de funciones para agregar y mostrar elementos, además de gestionar un menú interactivo que controle el flujo de la aplicación. Finalmente, se pretende verificar que el programa genere un ticket de salida preciso que incluya el desglose de los productos y la sumatoria total de los costos, validando así la correcta integración de la lógica algorítmica y el manejo de memoria durante la ejecución.